\documentclass[a4paper,12pt]{article}
\usepackage[T2A]{fontenc}
\usepackage[utf8]{inputenc}
\usepackage[russian]{babel}
\usepackage{amsmath,amssymb,mathtools}
\usepackage{geometry}
\geometry{left=19mm,right=19mm,top=20mm,bottom=21mm}
\usepackage{microtype}
\usepackage{booktabs,tabularx,array}
\usepackage{xcolor,tikz,hyperref,fancyhdr,parskip,enumitem,icomma}
\definecolor{accent}{HTML}{007579}
\definecolor{deep}{HTML}{253139}
\definecolor{err}{HTML}{A12D38}
\definecolor{muted}{HTML}{57656B}
\hypersetup{colorlinks=true,linkcolor=accent,urlcolor=accent,
 pdftitle={Ревью домашней работы. Комплексное повторение. 22 ноября 2026},
 pdfauthor={Лёвин Артём Александрович}}
\setlength{\parindent}{0pt}
\setlength{\parskip}{5pt plus 1pt}
\setlength{\headheight}{15pt}
\setlist{topsep=3pt,itemsep=2pt,parsep=0pt,leftmargin=22pt}
\pagestyle{fancy}
\fancyhf{}
\fancyhead[L]{\small\color{muted} Комплексное повторение}
\fancyhead[R]{\small\color{muted} 22 ноября 2026 г.}
\fancyfoot[C]{\small\thepage}
\renewcommand{\headrulewidth}{0.3pt}
\newcommand{\errlabel}{\textcolor{err}{\textbf{Ошибка.}}}
\newcommand{\keylabel}{\textcolor{accent}{\textbf{Ключевая идея.}}}
\newcommand{\resultlabel}{\textcolor{accent}{\textbf{Итог.}}}
\newcommand{\reviewtask}[2]{\clearpage\phantomsection\label{task:#1}\section*{\textcolor{accent}{№\,#1. #2}}}
\newcommand{\firstwrong}{\textbf{Первое место, требующее исправления.}}
\newcommand{\fix}{\textbf{Исправление.}}
\newcommand{\verify}{\textbf{Проверка.}}
\newcommand{\practice}{\medskip\textbf{Задача на закрепление.} }
\begin{document}
\thispagestyle{empty}
\begin{center}
\vspace*{31mm}
{\Large\color{accent} МАТЕМАТИКА}\par
\vspace{19mm}
{\huge\bfseries Ревью домашней работы}\par
\vspace{9mm}
{\Large\color{accent} Комплексное повторение}\par
\vspace{5mm}
{\large Степени, функции, движение, проценты, финансы, комбинаторика}\par
\vspace{19mm}
{\large Анастасия Павлова}\par
\vspace{16mm}
{\large 22 ноября 2026 г.}\par
\vfill
{\normalsize Лёвин Артём Александрович}\par
{\small эксклюзивно для Анастасии Павловой}\par
\vspace{11mm}
\begin{tikzpicture}
\draw[accent!75,line width=0.7pt] (-5.3,0) .. controls (-2.6,.65) and (-1.4,-.48) .. (0,0)
 .. controls (1.4,.48) and (2.6,-.65) .. (5.3,0);
\end{tikzpicture}
\vspace{3mm}
\end{center}

\clearpage
\section*{\textcolor{accent}{Сводная таблица замечаний}}
\small
\renewcommand{\arraystretch}{1.3}
\begin{tabularx}{\linewidth}{@{}>{\raggedright\arraybackslash}p{0.6cm}>{\raggedright\arraybackslash}p{2.4cm}>{\raggedright\arraybackslash}X >{\raggedright\arraybackslash}p{2.45cm}@{}}
\toprule
\textbf{№} & \textbf{Тема} & \textbf{Что исправить} & \textbf{Верный ответ}\\
\midrule
\hyperref[task:1]{\textbf{1}} & Степени & Корень переведён в неверную степень; неверная работа с показателями. & $x=\frac32$\\
\hyperref[task:2]{\textbf{2}} & Линейная функция & Вместо неизвестного $b$ записано $6$; решение не завершено. & $k=-2$, $b=3$\\
\hyperref[task:3]{\textbf{3}} & Векторы & Вместо решения оставлен прочерк. & $|AB|=10$, угол $90^\circ$\\
\hyperref[task:4]{\textbf{4}} & Катер & Неверно раскрыта разность квадратов; дальнейшая ошибка вычислений. & $9$ км/ч\\
\hyperref[task:5]{\textbf{5}} & Совместная работа & Время первого и второго мастеров перепутано. & Первый: $6$ ч; второй: $12$ ч\\
\hyperref[task:6]{\textbf{6}} & Растворы & Ошибка в массе соли; не найдена доля соли от всей смеси. & $28\%$\\
\hyperref[task:7]{\textbf{7}} & Проценты & Найдена начальная цена, но пропущено общее изменение. & $5000$ руб.; $+8\%$\\
\hyperref[task:8]{\textbf{8}} & Банковский вклад & Уравнение верно, вычисление ставки не закончено. & $10\%$ годовых\\
\hyperref[task:9]{\textbf{9}} & Погашение кредита & В выводе более выгодной названа не та схема. & Вторая дешевле на $1$ тыс. руб.\\
\hyperref[task:10]{\textbf{10}} & Разбиения & $6!$ считает перестановки; нужны разбиения на группы. & $10$; $P=\frac25$\\
\bottomrule
\end{tabularx}
\normalsize
\vspace{5mm}
По номеру в первом столбце можно перейти к подробному разбору. Формулы и выводы ниже проверены по условиям сборника «Комплексное повторение» от 16 ноября 2026 года.

\clearpage
\thispagestyle{empty}
\vspace*{\fill}
\begin{center}{\LARGE\bfseries\color{accent} Разбор ошибок}\end{center}
\vspace*{\fill}

\reviewtask{1}{Показательное уравнение}
\textbf{Условие.} Решите уравнение
\[
9^{x-1}=\sqrt{27}\cdot3^{x-2}.
\]
\textbf{Ваше решение.} На фотографии записан переход
\[
(3^2)^{x-1}=3^2\cdot3^{x-2},\qquad
2x-2=2(x-2),\qquad 2x-2=2x-4.
\]
\textbf{Ваш ответ.} Не получен.

\errlabel\ При замене корня степенью допущена ошибка; затем произведение степеней сведено к неверному показателю.

\firstwrong\ Левая часть $9^{x-1}=(3^2)^{x-1}$ записана верно. Первое неверное равенство возникает справа: $\sqrt{27}\ne3^2$. Действительно, $3^2=9$, а $\sqrt{27}=3\sqrt3$. Следующий переход $3^2\cdot3^{x-2}\to3^{2(x-2)}$ также неверен.

\textbf{Почему это неверно.} Корень квадратный равен степени с показателем $\frac12$:
\[
\sqrt{27}=\sqrt{3^3}=(3^3)^{1/2}=3^{3/2}.
\]
При умножении степеней с одинаковым основанием \emph{показатели складываются}:
\[
a^m\cdot a^n=a^{m+n}.
\]
Умножение показателей применяется к \emph{степени степени}: $(a^m)^n=a^{mn}$. Это два разных правила.

\keylabel\ Сначала все числа приводим к основанию $3$, отдельно преобразуя каждую степень. Затем складываем показатели множителей.

\fix\ От правильно записанной левой части переходим к
\[
3^{2(x-1)}=3^{3/2}\cdot3^{x-2}
=3^{\,3/2+x-2}.
\]
\textbf{Полное решение.} Поскольку основания равны $3>0$ и $3\ne1$, можно приравнять показатели:
\begin{align*}
2(x-1)&=\frac32+x-2,\\
2x-2&=x-\frac12,\\
x&=\frac32.
\end{align*}
\verify\ Слева при $x=\frac32$ имеем $9^{1/2}=3$. Справа
\[
\sqrt{27}\cdot3^{-1/2}
=\frac{3\sqrt3}{\sqrt3}=3.
\]
Обе части совпали.

\resultlabel\ $\boxed{x=\frac32}$.
\practice Решите уравнение $9^{x-2}=\sqrt3\cdot3^{x-1}$.

\reviewtask{2}{Функция по двум точкам}
\textbf{Условие.} Прямая $f(x)=kx+b$ проходит через точки $A(-2;7)$ и $B(2;-1)$. Найдите $k$, $b$, значение $f(-3)$ и пересечение с прямой $y=x+3$.

\textbf{Ваше решение.} Записано $f(x)=kx+6$ и затем $f(-3)=k(-3)+6$. Числовые ответы отсутствуют.

\textbf{Ваш ответ.} Не указан.

\errlabel\ Коэффициент $b$ был заменён числом $6$ без основания в условии, после чего решение остановилось.

\firstwrong\ Последний верный элемент --- общий вид линейной функции $f(x)=kx+b$. Первый ошибочный шаг --- запись $f(x)=kx+6$: значение свободного члена ещё неизвестно.

\textbf{Почему это неверно.} В выражении $kx+b$ число $k$ задаёт наклон прямой, а $b$ --- значение функции при $x=0$. Условие даёт две точки и позволяет найти оба числа из системы уравнений. Ни одна из координат данных точек не позволяет без вычислений присвоить $b$ значение $6$.

\keylabel\ Подставьте координаты обеих точек, решите систему для $k$ и $b$ и только затем вычисляйте значение функции и пересечение.

\fix\ Сохраняем $f(x)=kx+b$ и используем оба условия:
\[
\begin{cases}-2k+b=7,\\2k+b=-1.\end{cases}
\]
\textbf{Полное решение.} Вычтем первое уравнение из второго:
\[
(2k+b)-(-2k+b)=-1-7,
\quad4k=-8,\quad k=-2.
\]
Подставляем $k=-2$ в $2k+b=-1$:
\[
-4+b=-1,\qquad b=3.
\]
Следовательно, $f(x)=-2x+3$. Значение в заданной точке:
\[
f(-3)=-2\cdot(-3)+3=9.
\]
Точка пересечения с прямой $y=x+3$ удовлетворяет уравнению
\[
-2x+3=x+3,\quad -3x=0,\quad x=0,\quad y=3.
\]
\verify\ $f(-2)=4+3=7$ и $f(2)=-4+3=-1$; обе исходные точки принадлежат найденной прямой. При $x=0$ обе прямые имеют ординату $3$.

\resultlabel\ $\boxed{k=-2,\ b=3,\ f(-3)=9,\ (0;3)}$.
\practice Линейная функция $f(x)=kx+b$ проходит через $A(-1;5)$ и $B(3;-3)$. Найдите $k$, $b$, $f(-2)$ и точку пересечения с прямой $y=x+1$.

\reviewtask{3}{Векторы и угол между ними}
\textbf{Условие.} Даны $A(-2;1)$, $B(4;9)$, $C(6;-5)$. Найдите координаты $\overrightarrow{AB}$ и $\overrightarrow{AC}$, длину $AB$ и угол $BAC$.

\textbf{Ваше решение и ответ.} Рядом с номером записан прочерк; вычисления и ответ отсутствуют.

\errlabel\ Задание осталось без решения. Ошибочного вычислительного перехода здесь нет; требуется полностью построить решение.

\firstwrong\ Последний правильный шаг и первый неправильный шаг выделить невозможно, поскольку никаких математических действий для №\,3 не записано. Отправная точка --- координаты из условия.

\textbf{Что необходимо знать.} Чтобы найти координаты вектора из точки $A(x_A;y_A)$ в точку $B(x_B;y_B)$, вычитаем координаты начала из координат конца:
\[
\overrightarrow{AB}=(x_B-x_A;\ y_B-y_A).
\]
Длина вектора $(u;v)$ равна $\sqrt{u^2+v^2}$. Если скалярное произведение двух ненулевых векторов равно нулю, угол между ними прямой.

\keylabel\ Постройте оба вектора из общего начала $A$; для угла $BAC$ именно эти два вектора образуют нужные стороны.

\textbf{Полное решение.} Находим координаты:
\begin{align*}
\overrightarrow{AB}&=(4-(-2);\ 9-1)=(6;8),\\
\overrightarrow{AC}&=(6-(-2);\ -5-1)=(8;-6).
\end{align*}
Длина отрезка $AB$:
\[
AB=\sqrt{6^2+8^2}=\sqrt{36+64}=\sqrt{100}=10.
\]
Скалярное произведение векторов:
\[
\overrightarrow{AB}\cdot\overrightarrow{AC}
=6\cdot8+8\cdot(-6)=48-48=0.
\]
Векторы ненулевые, следовательно, они перпендикулярны и $\angle BAC=90^\circ$.

\verify\ $|\overrightarrow{AC}|=\sqrt{8^2+(-6)^2}=10$. Равные длины согласуются с полученным прямым углом, а нулевое скалярное произведение подтверждает перпендикулярность.

\resultlabel
\[\boxed{\overrightarrow{AB}=(6;8),\quad\overrightarrow{AC}=(8;-6)}\]
\[\boxed{AB=10,\qquad\angle BAC=90^\circ}.\]
\practice Даны $A(1;-2)$, $B(4;2)$, $C(5;-5)$. Найдите $\overrightarrow{AB}$, $\overrightarrow{AC}$, $AB$ и угол $BAC$.

\reviewtask{4}{Движение катера}
\textbf{Условие.} Катер проходит $24$ км против течения на $2$ часа дольше, чем по течению. Скорость течения $3$ км/ч. Найдите собственную скорость катера и укажите ограничение на неё.

\textbf{Ваше решение.} Верно составлено уравнение
\[
\frac{24}{x-3}-\frac{24}{x+3}=2,
\]
затем получено $24(x+3)-24(x-3)=2(x-3)(x+3)$. Далее в раскрытии скобок появляется выражение $x^2-3$ вместо $x^2-9$, и в результате записано $x=7$.

\textbf{Ваш ответ.} $7$ км/ч.

\errlabel\ Неверно раскрыто произведение сопряжённых скобок; затем ошибочно получено число под корнем.

\firstwrong\ Последний достоверно правильный шаг:
\[
24(x+3)-24(x-3)=2(x-3)(x+3).
\]
Первый неверный шаг --- замена $(x-3)(x+3)$ на $x^2-3$.

\textbf{Почему это неверно.} По формуле разности квадратов
\[
(a-b)(a+b)=a^2-b^2.
\]
Здесь $a=x$, $b=3$, следовательно, $(x-3)(x+3)=x^2-3^2=x^2-9$. В исходной записи потерян квадрат числа $3$. При дальнейших преобразованиях дополнительно нарушено вычисление квадратного корня: нельзя получить $x^2=49$ из $138=2x^2$.

\keylabel\ В задачах на движение по течению сначала фиксируем ограничение $x>3$, затем работаем с разностью квадратов и проверяем ответ подстановкой во времена.

\fix\ Возвращаемся к последнему верному равенству:
\[
24(x+3)-24(x-3)=2(x^2-9).
\]
\textbf{Полное решение.} Собственная скорость катера $x$ км/ч. Против течения скорость $x-3$, по течению --- $x+3$. Так как катер движется против течения, $x>3$. По формуле «время = расстояние : скорость»:
\[
\frac{24}{x-3}-\frac{24}{x+3}=2.
\]
Умножаем на $(x-3)(x+3)>0$:
\begin{align*}
24(x+3)-24(x-3)&=2(x^2-9),\\
24x+72-24x+72&=2x^2-18,\\
144&=2x^2-18,\\
162&=2x^2,\\
x^2&=81,\\
x&=\pm9.
\end{align*}
По ограничению $x>3$ оставляем только $x=9$.

\verify\ Против течения скорость $9-3=6$ км/ч, время $24:6=4$ ч. По течению скорость $9+3=12$ км/ч, время $24:12=2$ ч. Разница действительно $2$ часа.

\resultlabel\ $\boxed{9\text{ км/ч}}$, ограничение $\boxed{x>3}$.
\practice Катер проходит $24$ км против течения на $1$ час дольше, чем по течению. Скорость течения $1$ км/ч. Найдите собственную скорость катера и запишите ограничение.

\reviewtask{5}{Совместная работа мастеров}
\textbf{Условие.} Первый мастер выполняет заказ на $6$ часов быстрее второго; вместе они выполняют работу за $4$ часа. Найдите время каждого при самостоятельной работе.

\textbf{Ваше решение.} Составлено уравнение
\[
\frac1{x+6}+\frac1x=\frac14,
\]
получены корень $x=6$ и второе время $x+6=12$; в ответе записано: первый мастер --- $12$ часов, второй --- $6$ часов.

\textbf{Ваш ответ.} Первый: $12$ ч; второй: $6$ ч.

\errlabel\ В таблице неверно соотнесены переменная и исполнители. Из-за этого итоговое распределение времени противоречит условию.

\firstwrong\ Первое содержательно неверное место --- строка таблицы, где первому мастеру приписано время $x+6$, а второму --- $x$. Ведь $x+6>x$, то есть при такой записи первый мастер оказался бы медленнее.

\textbf{Почему это неверно.} Производительность обратно пропорциональна времени выполнения одного и того же заказа. Если первый работает быстрее, его время меньше на $6$ часов. Уравнение само по себе симметрично относительно слагаемых, поэтому вычисления $6$ и $12$ получились верными, но их отнесение к мастерам неверно.

\keylabel\ Перед составлением уравнения проверьте смысл обозначений: «быстрее на шесть часов» относится к меньшему времени, а не к большей производительности, выраженной в часах.

\fix\ Обозначим время первого мастера через $x$ ч, тогда время второго $x+6$ ч. Совместная производительность равна сумме индивидуальных.

\textbf{Полное решение.} За один час первый выполняет $\frac1x$ заказа, второй --- $\frac1{x+6}$. Вместе за час они выполняют $\frac14$ заказа:
\[
\frac1x+\frac1{x+6}=\frac14,\qquad x>0.
\]
Умножаем на $4x(x+6)$:
\begin{align*}
4(x+6)+4x&=x(x+6),\\
8x+24&=x^2+6x,\\
x^2-2x-24&=0,\\
(x-6)(x+4)&=0.
\end{align*}
Корень $x=-4$ не подходит по смыслу, следовательно, $x=6$. Второй мастер работает $6+6=12$ часов.

\verify\ За час они выполняют
\[
\frac16+\frac1{12}=\frac2{12}+\frac1{12}=\frac14
\]
заказа. Вместе работа займёт $4$ часа; первый выполняет её быстрее второго ровно на $6$ часов.

\resultlabel\ $\boxed{\text{Первый: }6\text{ ч; второй: }12\text{ ч}}$.
\practice Первый мастер выполняет заказ на $5$ часов быстрее второго. Работая вместе, они выполняют заказ за $6$ часов. Найдите время каждого мастера при самостоятельной работе.

\reviewtask{6}{Испарение и смешивание растворов}
\textbf{Условие.} Было $8$ кг $30\%$-го раствора соли. Испарили $2$ кг воды, добавили $4$ кг $10\%$-го раствора. Найдите массовую долю соли в полученной смеси.

\textbf{Ваше решение.} В таблице записана масса соли $0{,}24$ кг для $8$ кг $30\%$-го раствора, затем $0{,}4$ кг из добавленного раствора. Итог: $0{,}64$.

\textbf{Ваш ответ.} $0{,}64$ (без указания процентов).

\errlabel\ Масса соли первоначально посчитана с ошибкой в десятичном разряде. Кроме того, конечная доля соли должна делиться на полную массу смеси.

\firstwrong\ Первая ошибка возникает в произведении $8\cdot0{,}3$: оно равно $2{,}4$, а не $0{,}24$. Значение $0{,}4=4\cdot0{,}1$ для добавленного раствора вычислено верно.

\textbf{Почему это неверно.} Тридцать процентов --- это $\frac{30}{100}=0{,}3$. Поэтому тридцать процентов от восьми килограммов --- $2{,}4$ килограмма. При испарении чистой воды соль остаётся в сосуде; изменяется только общая масса раствора. После смешивания процент находят как отношение массы соли к массе \emph{всей} новой смеси.

\keylabel\ В задачах на растворы ведите две отдельные величины: массу смеси и массу растворённого вещества. Испарение воды уменьшает только первую.

\fix\ Заменяем неверную массу соли $0{,}24$ кг на $2{,}4$ кг, добавляем $0{,}4$ кг соли и делим результат на массу полученной смеси.

\textbf{Полное решение.}
\begin{enumerate}
\item Исходная масса соли: $8\cdot0{,}30=2{,}4$ кг.
\item После испарения $2$ кг воды соли остаётся $2{,}4$ кг, а масса раствора становится $8-2=6$ кг.
\item С добавленными $4$ кг $10\%$-го раствора поступает $4\cdot0{,}10=0{,}4$ кг соли.
\item Итоговая масса соли $2{,}4+0{,}4=2{,}8$ кг, масса смеси $6+4=10$ кг.
\end{enumerate}
Следовательно,
\[
w=\frac{2{,}8}{10}=0{,}28=28\%.
\]
\verify\ В десяти килограммах смеси соль занимает $2{,}8$ кг, вода $10-2{,}8=7{,}2$ кг; вместе $10$ кг. Значение $28\%$ находится между концентрациями раствора после выпаривания и добавленного слабого раствора.

\resultlabel\ $\boxed{28\%}$.
\practice В сосуде $10$ кг $20\%$-го раствора соли. Испарили $3$ кг воды и добавили $5$ кг $40\%$-го раствора той же соли. Найдите массовую долю соли в смеси (в процентах).

\reviewtask{7}{Последовательные изменения цены}
\textbf{Условие.} Цена выросла на $20\%$, затем уменьшилась на $10\%$. Итоговая цена $5400$ руб. Найдите первоначальную цену и полное изменение в процентах от первоначальной.

\textbf{Ваше решение.} По таблице: первоначальная цена $5000$ руб., после увеличения $6000$ руб., после уменьшения $5400$ руб.

\textbf{Ваш ответ.} $5000$ руб.; процент полного изменения не указан.

\errlabel\ Один из двух требуемых результатов отсутствует. Рассчитана цена, но не рассчитано изменение в процентах.

\firstwrong\ Действия $5000\cdot1{,}2=6000$ и $6000\cdot0{,}9=5400$ правильны. Ошибка возникает при завершении ответа после числа $5000$: в условии отдельно требуется определить общий процент изменения цены.

\textbf{Почему решение неполное.} Проценты второго изменения считаются от новой, уже увеличенной цены. Итоговая стоимость выше первоначальной на $5400-5000=400$ руб. Чтобы найти процент, нужно отнести $400$ именно к первоначальной цене $5000$ руб.

\keylabel\ Для последовательных процентных изменений используйте общий множитель: увеличение на $p\%$ соответствует $1+p/100$, уменьшение на $q\%$ соответствует $1-q/100$.

\fix\ К имеющемуся верному ответу $5000$ руб. добавляем вычисление процента:
\[
\frac{5400-5000}{5000}\cdot100\%=8\%.
\]
\textbf{Полное решение.} Пусть $x$ --- цена до изменений. Тогда
\[
1{,}2\cdot0{,}9\cdot x=5400,\qquad1{,}08x=5400,
\]
\[
x=\frac{5400}{1{,}08}=5000\text{ руб.}
\]
Общий коэффициент изменения $1{,}08$ означает повышение на $0{,}08$ первоначальной цены, то есть на $8\%$.

\verify\ $5000\to6000\to5400$; итоговый прирост $400$ руб. Доля $400:5000=0{,}08$.

\resultlabel\ $\boxed{5000\text{ руб.; повышение на }8\%}$.
\practice Цену товара сначала повысили на $15\%$, затем понизили на $10\%$. После изменений товар стоит $5175$ руб. Найдите начальную цену и общий процент изменения.

\reviewtask{8}{Вклад с капитализацией}
\textbf{Условие.} В начале первого года внесли $600$ тыс. руб. В конце первого года после начисления процентов внесли ещё $120$ тыс. руб. В конце второго года на вкладе $858$ тыс. руб. Ставка одинакова. Найдите годовую ставку.

\textbf{Ваше решение.} Верно составлено $600k^2+120k=858$ и затем $100k^2+20k-143=0$. Для него найден дискриминант $D=57600$, после чего решение обрывается на $\sqrt{57600}$.

\textbf{Ваш ответ.} Не получен.

\errlabel\ Решение остановилось перед нахождением корня и переводом коэффициента роста в проценты.

\firstwrong\ Все приведённые основные уравнения и число $D=57600$ правильны. Первый недостающий шаг --- вычислить $\sqrt D=240$, затем получить допустимый коэффициент $k$.

\textbf{Почему этого недостаточно.} Переменная $k$ --- коэффициент годового роста вклада, поэтому она сама по себе ещё не является процентной ставкой. Для положительной ставки требуется $k>1$, а процент вычисляется по формуле $p=(k-1)\cdot100\%$.

\keylabel\ При капитализации проценты начисляются на текущую сумму: сумма $600$ проходит через два года, пополнение $120$ --- только через второй год.

\fix\ Продолжаем от Вашего правильного дискриминанта:
\[
\sqrt{57600}=240,
\quad k=\frac{-20\pm240}{200}.
\]
\textbf{Полное решение.} Пусть $k=1+p/100$ --- годовой коэффициент. После первого начисления имеется $600k$, после пополнения --- $600k+120$. После второго года:
\[
(600k+120)k=858.
\]
Раскрываем скобки, приводим уравнение:
\begin{align*}
600k^2+120k-858&=0,\\
100k^2+20k-143&=0.
\end{align*}
Дискриминант
\[
D=20^2-4\cdot100\cdot(-143)=400+57200=57600.
\]
Корни:
\[
k_1=\frac{-20+240}{200}=1{,}1,
\qquad k_2=\frac{-20-240}{200}=-1{,}3.
\]
Отрицательный коэффициент роста не соответствует условию вклада. Следовательно,
\[
p=(1{,}1-1)\cdot100\%=10\%.
\]
\verify\ $600\cdot1{,}1=660$ тыс. руб.; после пополнения $780$ тыс. руб.; через год $780\cdot1{,}1=858$ тыс. руб.

\resultlabel\ $\boxed{10\%\text{ годовых}}$.
\practice В начале первого года внесли $400$ тыс. руб. В конце первого года после начисления процентов пополнили вклад на $120$ тыс. руб. В конце второго года стало $616$ тыс. руб. Определите постоянную годовую ставку, если проценты капитализировались ежегодно.

\reviewtask{9}{Две схемы погашения кредита}
\textbf{Условие.} Кредит $110$ тыс. руб. на два года под $20\%$ годовых. Сравните два равных ежегодных платежа и ежегодное погашение половины основного долга плюс проценты на остаток. Найдите платежи, переплату и более выгодный вариант.

\textbf{Ваше решение.} Получены равные платежи по $72$ тыс. руб. и переплата $34$ тыс. руб.; при второй схеме платежи $77$ и $66$ тыс. руб., переплата $33$ тыс. руб. Разница $34-33=1$ тыс. руб. В заключении более выгодной названа \emph{первая} схема.

\textbf{Ваш ответ.} «Первая схема более выгодная, чем вторая» (по видимой заключительной строке).

\errlabel\ Итоговый словесный вывод противоречит Вашим же верным числовым результатам.

\firstwrong\ До сравнения переплат всё вычислено согласованно: $34$ тыс. руб. по первой схеме и $33$ тыс. руб. по второй. Первое ошибочное утверждение --- назвать более выгодной первую схему, хотя $34>33$.

\textbf{Почему это неверно.} При одинаковой выданной сумме кредита более выгодна для заёмщика схема с \emph{меньшей суммарной переплатой}. Разница в $1$ тыс. руб. говорит в пользу второй схемы.

\keylabel\ Сравнивайте общий итог платежей или переплату, используя один критерий и одинаковые единицы измерения. Большая переплата означает более дорогой кредит.

\fix\ Сохраняем Ваши верные числовые расчёты, меняем только заключение: вторая схема дешевле на $1$ тыс. руб.

\textbf{Полное решение.} Все суммы ниже измеряются в тысячах рублей.

\emph{Схема а: два одинаковых платежа $A$.} После первого года долг с процентами $110\cdot1{,}2=132$. После платежа $A$ остаток $132-A$. Перед вторым платежом он вырастает до $(132-A)\cdot1{,}2$. Поэтому
\[
(132-A)\cdot1{,}2=A.
\]
Отсюда $158{,}4=2{,}2A$, то есть $A=72$. Всего выплачено $144$, переплата $144-110=34$.

\emph{Схема б: каждый год погашается половина основного долга.} Первая половина долга $55$, проценты первого года $110\cdot0{,}2=22$, значит, первый платёж $55+22=77$. Остаток долга $55$; проценты второго года $55\cdot0{,}2=11$. Второй платёж $55+11=66$. Всего выплачено $143$, переплата $143-110=33$.

Сравним: $34-33=1$ тыс. руб.; дешевле схема б.

\verify\ В схеме а после первого платежа остаток $132-72=60$, во втором году $60\cdot1{,}2=72$, значит, остаток закрыт полностью. В схеме б основной долг погашается двумя частями по $55$.

\resultlabel\ $\boxed{\text{Вторая схема дешевле на }1\text{ тыс. руб.}}$.
\practice Кредит $220$ тыс. руб. выдан на $2$ года под $20\%$ годовых. Сравните два равных ежегодных платежа и погашение половины первоначального долга ежегодно с процентами на остаток. Найдите платежи, переплату и разницу в стоимости схем.

\reviewtask{10}{Разбиения и вероятность}
\textbf{Условие.} Шесть участников случайно распределяют на две \emph{неименованные} группы по три человека; все разбиения равновероятны. Сколько существует разбиений? Какова вероятность, что Аня и Борис будут в одной группе?

\textbf{Ваше решение.} Записано $6!=720$. Подсчёта разбиений и вероятности нет.

\textbf{Ваш ответ.} Указано $720$; вероятность не определена.

\errlabel\ Посчитано количество упорядоченных расположений шести людей, тогда как условие требует количество разбиений на две неименованные тройки.

\firstwrong\ Сначала правильно выделено, что участников шесть. Первый неверный переход --- использование $6!$ в качестве количества разбиений. $6!$ отвечает на задачу «Сколькими способами можно выстроить шесть разных людей в ряд?».

\textbf{Почему это неверно.} В пределах каждой тройки порядок людей роли не играет, а две группы не имеют названий «первая» и «вторая». Поэтому один и тот же состав групп многократно появляется среди $720$ перестановок. Нужно убрать внутренний порядок и перестановку самих групп:
\[
\frac{6!}{3!\,3!\,2!}=10.
\]
Фактор $3!$ для каждой группы убирает внутренние перестановки, а $2!$ --- обмен двух неименованных групп.

\keylabel\ Перед подсчётом определите, что считать разными результатами: порядок участников внутри группы и порядок самих групп здесь значения не имеют.

\fix\ Вместо $6!$ считаем способы выбрать троих из шести и делим на два, потому что выбор одной из двух неименованных групп повторяет то же разбиение при выборе другой:
\[
N=\frac{\binom63}{2}=\frac{20}{2}=10.
\]
\textbf{Полное решение.} Всего $10$ равновозможных разбиений. Найдём благоприятные. Зафиксируем, что в одной тройке находятся Аня и Борис. Третьего участника их группы можно выбрать из четырёх остальных: $4$ способа. Оставшаяся тройка при этом полностью определена, и каждое разбиение посчитано ровно один раз. Значит,
\[
N_{\text{благ}}=4,\qquad P=\frac4{10}=\frac25=0{,}4.
\]

\verify\ Можно рассуждать иначе: зафиксируем группу Ани. Из пяти остальных участников в её группу случайно попадут два. Вероятность, что среди двух окажется Борис, равна $\frac25$; это совпадает с результатом подсчёта разбиений.

\resultlabel\ $\boxed{10\text{ разбиений};\quad P=\frac25=0{,}4}$.
\practice Восемь участников делят случайно на две неименованные группы по четыре человека; все разбиения равновероятны. Найдите число разбиений и вероятность того, что два определённых человека окажутся в одной группе.

\vfill
\begin{center}\small\color{muted} Лёвин Артём Александрович --- эксклюзивно для Анастасии Павловой\end{center}
\end{document}